\subsubsection{Matrix Product Representation and Tensorized MOM Compilation of Fiber Multiplicities: Non-Commutative Counting under Finite-Window Inverse Codes}\label{subsubsec:pom-multiplicity-matrix-product-tensor-mom}
In the main text's collision spectrum framework, \(S_q(m)=\sum_{x\in X_m}d_m(x)^q\) is treated as a ``\(q\)-fold fiber product'' counting object (Definition \ref{def:pom-moment-q}), unified through finite-state moment-kernels into the Perron--Frobenius/\(\zeta\) interface (Proposition \ref{prop:pom-moment-kernel-exists} and Definition \ref{def:pom-moment-pressure}).
This paragraph supplements a finer structural layer: under the regime where a \emph{finite-window inverse code} exists, the individual fiber volume \(d_m(x)\) itself admits a matrix product counting formula; moreover, for any integer order \(q\), the collision kernel of \((\mathrm{MOM}_q)\) can be explicitly written as the Kronecker coupling of ``the syntactic two-state golden skeleton \(\times\) the \(q\)-fold synchronous direct product of the inverse-code transfer.''

\paragraph{Finite-window inverse code \texorpdfstring{$\Rightarrow$}{=>} matrix product counting of fiber volumes}
\begin{theorem}[Matrix product representation of fiber multiplicities and minimal ambiguity-state dimension]\label{thm:pom-fiber-multiplicity-matrix-product}
Let \(\Fold_\infty:\Omega_\infty=\{0,1\}^{\NN}\to X_\infty\) be the binary folding factor, and let \(X_\infty\subseteq\{0,1\}^{\NN}\) be the golden mean shift (forbidding adjacent \(11\)).
Suppose \(\Fold_\infty\) admits a finite-window inverse code \(\Psi\), so that ``given the output prefix \(x_1\cdots x_m\),'' all legitimate preimages can be updated and counted position-by-position via a finite-state deterministic system (equivalently, one can take a follower-separated right-resolving presentation and perform counting readout on it).
Then there exist a finite state set \(\mathcal S\), boundary vectors \(u,v\in\NN^{\mathcal S}\), and two matrices
\[
T_0,T_1\in M_{|\mathcal S|}(\NN),
\]
such that for any legitimate macroscopic word \(x=x_1\cdots x_m\in X_m\) (\(m\ge 1\)) the exact counting formula holds:
\[
\boxed{\;
d_m(x)=u^{\mathsf T}T_{x_1}T_{x_2}\cdots T_{x_m}v\; }.
\]
Furthermore, among all choices satisfying the above representation, the minimum possible dimension \(|\mathcal S|\) equals the cardinality of the finite complete set of right ambiguity classes (follower ambiguity classes);
therefore, when \(\Psi\) is taken to be the canonical inverse code induced by the minimal online delay/minimal synchronization structure, this minimum dimension is uniquely determined by \(\Fold_\infty\).
\end{theorem}
\begin{proof}
By hypothesis, the inverse code can be realized as a deterministic finite-state counting automaton: its states record ``the ambiguity classes still possible given the output prefix,'' and upon reading the next output symbol \(b\in\{0,1\}\), a deterministic update is performed.
Let \(\mathcal S\) be the (reachable) state set of this counting automaton, and let \((T_b)_{b\in\{0,1\}}\) be the transition counting matrices decomposed by output symbol (the counting semantics at each step are given by the branching numbers produced by the inverse code, so entries are nonneg\-ative integers).
Take \(u\) as the initial state distribution (determined by boundary/starting admissible states), and \(v\) as the terminal readout vector (summing over admissible terminating states). Then the total number of preimages of the length-\(m\) output word \(x_1\cdots x_m\) is precisely the path count along this deterministic counting automaton, yielding the matrix product formula.
\par
The minimality conclusion is standard minimization of deterministic finite automata: merging states by the equivalence relation ``having the same follower set (identical right extension family)'' yields the minimal right-resolving follower-separated representation, whose number of equivalence classes is the stated right ambiguity class complete set cardinality; any matrix product representation must distinguish these equivalence classes, so its dimension cannot be less than this cardinality, while the minimization construction achieves this lower bound.
\end{proof}

\paragraph{Tensorized MOM compilation: golden syntactic skeleton \texorpdfstring{$\otimes$}{x} \(q\)-fold inverse-code synchronous direct product}
\begin{proposition}[Tensor-decomposition-type collision kernel for MOM\texorpdfstring{$\,_q$}{\_q}]\label{prop:pom-momq-tensor-kernel-construction}
Under the setting of Theorem \ref{thm:pom-fiber-multiplicity-matrix-product}, let
\[
G_0:=\begin{pmatrix}1&0\\[2pt]1&0\end{pmatrix},
\qquad
G_1:=\begin{pmatrix}0&1\\[2pt]0&0\end{pmatrix}
\in M_2(\{0,1\})
\]
be the symbol-decomposed transition matrices of the golden mean two-state skeleton (satisfying \(G_0+G_1=\bigl(\begin{smallmatrix}1&1\\1&0\end{smallmatrix}\bigr)\)).
For any integer \(q\ge 1\), define the explicit \(q\)-th order collision transfer kernel
\[
A_q^{\mathrm{ten}}
\;:=\;
G_0\otimes T_0^{\otimes q}\;+\;G_1\otimes T_1^{\otimes q}
\ \in\ M_{2|\mathcal S|^q}(\NN).
\]
Let \(e_0:=(1,0)^{\mathsf T}\in\NN^2\), \(\mathbf 1:=(1,1)^{\mathsf T}\in\NN^2\), and take the boundary vectors
\[
\alpha_q:=e_0\otimes u^{\otimes q},
\qquad
\beta_q:=\mathbf 1\otimes v^{\otimes q}.
\]
Then for all \(m\ge 1\) the linear recurrence counting representation holds:
\[
\boxed{\;
S_q(m)=\sum_{x\in X_m}d_m(x)^q
\;=\;
\alpha_q^{\mathsf T}\bigl(A_q^{\mathrm{ten}}\bigr)^{m}\beta_q\; }.
\]
Therefore, the ``finite-state compilation'' of \((\mathrm{MOM}_q)\) can be taken as a tensor-decomposition-type structure: the \(q\)-body overlap constraint at the kernel level is equivalent to the \(q\)-fold synchronous direct product of the inverse-code transfer, Kronecker-coupled with the golden syntactic two-state skeleton.
\end{proposition}
\begin{proof}
By Theorem \ref{thm:pom-fiber-multiplicity-matrix-product},
\(
d_m(x)=u^{\mathsf T}T_{x_1}\cdots T_{x_m}v
\).
Hence
\[
d_m(x)^q
=
\bigl(u^{\otimes q}\bigr)^{\mathsf T}\,
\bigl(T_{x_1}^{\otimes q}\bigr)\cdots
\bigl(T_{x_m}^{\otimes q}\bigr)\,
v^{\otimes q}.
\]
Summing over all legitimate words \(x\in X_m\) can be realized through the path expansion of the golden mean two-state skeleton: at each step one chooses a symbol \(b\in\{0,1\}\) and updates the syntactic state via \(G_b\), while simultaneously applying the synchronous update \(T_b^{\otimes q}\) on \(\mathcal S^{\times q}\).
Therefore the weighted path sum is exactly \(\alpha_q^{\mathsf T}(A_q^{\mathrm{ten}})^m\beta_q\). \qed
\end{proof}

\paragraph{Disjoint tensor base: explicit spectrum of \texorpdfstring{$M^{\otimes q}$}{M^⊗q} and the fourth-order integer critical point}
Let
\(
M:=G_0+G_1=\bigl(\begin{smallmatrix}1&1\\[2pt]1&0\end{smallmatrix}\bigr)
\).
This two-state matrix not only encodes the golden mean syntactic skeleton, but also provides a more primitive tensor hardcore object: the adjacency operator of the disjoint graph.

\begin{proposition}[Tensor power characterization of the disjoint graph and the universal spectral base \texorpdfstring{$\QQ(\sqrt5)$}{Q(sqrt5)}]\label{prop:pom-disjointness-graph-tensor-power-spectrum}
Let \(\mathcal{D}_q\) be the disjoint graph on vertex set \(2^{[q]}\): for \(S,T\subseteq[q]\), let \(S\sim T\) if and only if \(S\cap T=\varnothing\).
In the natural basis (labeled by \(\{0,1\}^q\cong 2^{[q]}\)), the adjacency operator of \(\mathcal{D}_q\) is exactly
\[
\boxed{\;A(\mathcal{D}_q)=M^{\otimes q}\;}
\qquad\Bigl(M=\begin{pmatrix}1&1\\[2pt]1&0\end{pmatrix}\Bigr).
\]
Hence its spectrum is completely explicit: if \(\varphi=(1+\sqrt5)/2\), then all eigenvalues are
\[
(-1)^k\varphi^{\,q-2k}\qquad(0\le k\le q),
\]
with algebraic multiplicity \(\binom{q}{k}\).
Therefore the spectral decomposition of the disjoint tensor base is governed by \(\QQ(\sqrt5)\) at each \(q\).
\end{proposition}
\begin{proof}
Identifying \(S\subseteq[q]\) with the indicator vector \(s\in\{0,1\}^q\), \(S\cap T=\varnothing\) is equivalent to: for each coordinate \(i\), the local pair \((s_i,t_i)=(1,1)\) is not allowed.
Hence the adjacency matrix decomposes by coordinates as the Kronecker product of \(q\) independent \(2\times 2\) local admissibility matrices, i.e., \(A(\mathcal{D}_q)=M^{\otimes q}\).
Since \(M\) has eigenvalues \(\varphi\) and \(-\varphi^{-1}\), the eigenvalues of the tensor power are all products \(\varphi^{q-k}(-\varphi^{-1})^k=(-1)^k\varphi^{q-2k}\), and the multiplicity is given by the binomial expansion. \qed
\end{proof}

\begin{corollary}[\texorpdfstring{$q=4$}{q=4} as the unique integer critical point of the Ramanujan ratio at the disjoint graph level]\label{cor:pom-disjointness-ramanujan-q4-critical}
Let \(r_q^{\mathrm{disj}}:=\rho(A(\mathcal{D}_q))\) be the Perron root, and let \(\Lambda_q^{\mathrm{disj}}\) be the maximum-modulus non-Perron eigenvalue.
Then
\[
r_q^{\mathrm{disj}}=\varphi^{\,q},\qquad \Lambda_q^{\mathrm{disj}}=\varphi^{\,q-2},\qquad
\frac{\Lambda_q^{\mathrm{disj}}}{\sqrt{r_q^{\mathrm{disj}}}}=\varphi^{\,q/2-2}.
\]
Therefore for integer \(q\le 3\), \(\Lambda_q^{\mathrm{disj}}<\sqrt{r_q^{\mathrm{disj}}}\); for \(q=4\), the critical equation \(\Lambda_4^{\mathrm{disj}}=\sqrt{r_4^{\mathrm{disj}}}\) holds; and for integer \(q\ge 5\), necessarily \(\Lambda_q^{\mathrm{disj}}>\sqrt{r_q^{\mathrm{disj}}}\).
\end{corollary}
\begin{proof}
By Proposition \ref{prop:pom-disjointness-graph-tensor-power-spectrum}, \(r_q^{\mathrm{disj}}=\varphi^q\), and the maximum non-Perron modulus comes from the \(k=1\) branch, hence \(\Lambda_q^{\mathrm{disj}}=\varphi^{q-2}\).
Substituting yields the closed-form ratio, and the uniqueness of the critical point follows from \(\varphi>1\). \qed
\end{proof}

\begin{conclusion}[Quantized critical point of the Kronecker disjoint kernel: the critical \(q\) for \texorpdfstring{$S_q$}{Sq}-equivariant compression is quantized to \(4k_0\)]\label{con:pom-disjointness-sq-equivariant-critical-quantization}
Following Proposition \ref{prop:pom-disjointness-graph-tensor-power-spectrum}, let
\(
M=\bigl(\begin{smallmatrix}1&1\\[2pt]1&0\end{smallmatrix}\bigr)
\)
and consider the action of \(M^{\otimes q}\) on \((\RR^2)^{\otimes q}\).
Let \(W\subset(\RR^2)^{\otimes q}\) be any \(S_q\)-invariant subspace containing the Perron mode (equivalently, containing the Perron eigenvector of \(M^{\otimes q}\)).
Let
\[
T:=\mathrm{Proj}_W(M^{\otimes q})\mathrm{Proj}_W
\]
be the corresponding \(S_q\)-equivariant compression operator, and let \(k_0\) be the minimal nonzero ``weight order'' contained in \(W\): the smallest \(k\ge 1\) such that the projection of \(W\) onto the eigenspace of eigenvalue \((-1)^k\varphi^{q-2k}\) is nonzero.
Then the compression operator still satisfies \(r=\rho(T)=\varphi^q\), and its maximum non-Perron modulus is precisely
\[
\Lambda=\varphi^{q-2k_0},
\qquad
\frac{\Lambda}{\sqrt r}=\varphi^{\,q/2-2k_0}.
\]
Therefore the Ramanujan critical condition \(\Lambda\le \sqrt r\) is equivalent to \(q\le 4k_0\), and the critical point is quantized to
\[
\boxed{\;q_c=4k_0\;}
\]
(on integer \(q\), this manifests as a discrete threshold sequence).
In particular, \(k_0=1\) corresponding to the minimal nontrivial compression yields the criticality of \(q=4\).
\end{conclusion}

\begin{proposition}[\texorpdfstring{$S_q$}{Sq}-symmetric compression: closed-form \texorpdfstring{$(q+1)$}{q+1}-dimensional matrix and the spectral invariant of the fourth-order integer \(7\)]\label{prop:pom-disjointness-symmetric-compression-bq}
Let \(B_q\) be the operator induced by \(\mathcal{D}_q\) on the \(S_q\) orbit space (stratified by subset size).
Then \(B_q\) is a \((q+1)\times(q+1)\) matrix; in the basis labeled by ``size layers,'' its entries satisfy the explicit closed form
\[
\boxed{\;B_q(i,j)=\binom{q-j}{i}\;}\qquad(0\le i,j\le q).
\]
In particular, when \(q=4\), its characteristic polynomial factors strictly as
\[
(\lambda-1)(\lambda^2-7\lambda+1)(\lambda^2+3\lambda+1),
\]
with five eigenvalues
\(
1
\), \(
\frac{7\pm 3\sqrt5}{2}=\varphi^{\pm 4}
\), \(
\frac{-3\pm\sqrt5}{2}=-\varphi^{\pm 2}
\).
Therefore the ``coefficient \(7\)'' at the fourth-order level already appears as a spectral invariant in the minimal nontrivial symmetric compression of the disjoint tensor base.
\end{proposition}
\begin{proof}
For fixed \(T\subseteq[q]\) with \(|T|=j\), the number of \(i\)-element subsets disjoint from \(T\) is exactly the count of choosing \(i\)-element subsets from the complement \([q]\setminus T\) (of size \(q-j\)), namely \(\binom{q-j}{i}\), hence \(B_q(i,j)=\binom{q-j}{i}\).
When \(q=4\), directly computing and factoring the characteristic polynomial of this \(5\times 5\) matrix yields the stated closed form. \qed
\end{proof}

\paragraph{Lucas square spectral decomposition and closed-form trace formula: completely computable periodic dynamics of \texorpdfstring{$B_q$}{Bq}}
Let
\(
M=\bigl(\begin{smallmatrix}1&1\\[2pt]1&0\end{smallmatrix}\bigr)
\)
and denote the Lucas sequence by \(L_0=2,L_1=1\) with \(L_{n+1}=L_n+L_{n-1}\).

\begin{proposition}[\texorpdfstring{$B_q=\mathrm{Sym}^q(M)$}{Bq=Sym^q(M)} and explicit spectrum]\label{prop:pom-bq-is-symq-and-spectrum}
For any integer \(q\ge 1\), the symmetric compression kernel \(B_q\) is isomorphic to the \(q\)-th symmetric power of \(M\):
\[
\boxed{\;B_q\ \cong\ \mathrm{Sym}^q(M)\in \mathrm{GL}_{q+1}(\ZZ)\;}
\]
Hence all its eigenvalues are precisely
\[
\boxed{\;\mu_k=(-1)^k\varphi^{\,q-2k}\qquad(0\le k\le q)\;}
\]
and each mode has algebraic multiplicity \(1\) on \(B_q\).
\end{proposition}
\begin{proof}
Let \(V=\RR^2\) and take basis vectors \(e_0,e_1\).
Using the standard isomorphism, represent \(\mathrm{Sym}^q(V)\) as the space of homogeneous degree-\(q\) binary polynomials: formal variables \(x,y\) respectively correspond to \(e_0,e_1\), so \(\{x^{q-j}y^j\}_{j=0}^{q}\) forms a basis.
The matrix \(M\) in basis \((e_0,e_1)\) satisfies \(Me_0=e_0+e_1\), \(Me_1=e_0\), so it induces the substitution action
\((x,y)\mapsto (x+y,x)\) in this polynomial model.
Hence the image of \(x^{q-j}y^j\) is \((x+y)^{q-j}x^{j}\), whose expansion coefficients in the basis \(\{x^{q-i}y^i\}\) are precisely \(\binom{q-j}{i}\);
thus the matrix entries of \(\mathrm{Sym}^q(M)\) coincide with the closed form of Proposition \ref{prop:pom-disjointness-symmetric-compression-bq}, yielding \(B_q\cong\mathrm{Sym}^q(M)\).
\par
Since the eigenvalues of \(M\) are \(\varphi\) and \(-\varphi^{-1}\), the spectrum of the symmetric power consists of all products \(\varphi^{q-k}(-\varphi^{-1})^{k}=(-1)^k\varphi^{q-2k}\ (0\le k\le q)\), and the degree is \(q+1\), so all multiplicities are \(1\). \qed
\end{proof}

\begin{theorem}[Lucas square spectral decomposition: complete factorization of \texorpdfstring{$\det(I-uB_q^2)$}{det(I-uBq^2)}]\label{thm:pom-bq-lucas-square-spectrum-factorization}
For any integer \(q\ge 1\), the strict identity holds:
\[
\boxed{\;
\det(I-uB_q^2)
=(1-u)^{\mathbf 1_{\{q\ \mathrm{even}\}}}
\prod_{\substack{1\le d\le q\\ d\equiv q\ (\mathrm{mod}\ 2)}}
\bigl(u^2-L_{2d}\,u+1\bigr)\; } .
\]
In particular, when \(q=4\),
\[
\det(I-uB_4^2)=(1-u)(u^2-7u+1)(u^2-47u+1),
\]
where \(7=L_4\), \(47=L_8\) provide the integer-coefficient origins of the fourth-order square spectral trace.
\end{theorem}
\begin{proof}
By Proposition \ref{prop:pom-bq-is-symq-and-spectrum}, the eigenvalues of \(B_q\) are \(\mu_k=(-1)^k\varphi^{q-2k}\), so the eigenvalues of \(B_q^2\) are \(\lambda_k:=\mu_k^2=\varphi^{2(q-2k)}\).
When \(q\) is even, there exists a unique self-inverse mode \(k=q/2\) such that \(\lambda_{q/2}=1\), contributing the factor \(1-u\); the remaining eigenvalues pair as \(\lambda\leftrightarrow 1/\lambda\).
For any positive integer \(d\equiv q\ (\mathrm{mod}\ 2)\) (corresponding to \(\lambda=\varphi^{2d}\)), the Lucas even-index identity
\(
L_{2d}=\varphi^{2d}+\varphi^{-2d}
\)
holds; hence
\[
(1-u\varphi^{2d})(1-u\varphi^{-2d})=u^2-(\varphi^{2d}+\varphi^{-2d})u+1=u^2-L_{2d}u+1.
\]
Multiplying all reciprocal pair factors yields the conclusion. \qed
\end{proof}

\begin{theorem}[Closed-form trace formula: Lucas expansion of \texorpdfstring{$\mathrm{Tr}(B_q^n)$}{Tr(Bq^n)}]\label{thm:pom-bq-trace-lucas}
For any integers \(q\ge 1\), \(n\ge 1\), the strict closed form holds:
\[
\boxed{\;
\mathrm{Tr}(B_q^{\,n})
=
\sum_{k=0}^{\lfloor q/2\rfloor-\mathbf 1_{\{q\ \mathrm{even}\}}}
(-1)^{kn}\,L_{n(q-2k)}
\;+\;
\mathbf 1_{\{q\ \mathrm{even}\}}\cdot (-1)^{(q/2)n}\; } .
\]
Therefore the primitive orbit count for \(B_q\) (Möbius inversion)
\[
p_n(B_q):=\frac{1}{n}\sum_{d\mid n}\mu(d)\,\mathrm{Tr}\!\bigl(B_q^{\,n/d}\bigr)
\]
can be reduced in finitely many steps to a finite linear combination of Lucas numbers.
\end{theorem}
\begin{proof}
By Proposition \ref{prop:pom-bq-is-symq-and-spectrum},
\(
\mathrm{Tr}(B_q^{\,n})=\sum_{k=0}^{q}(-1)^{kn}\varphi^{n(q-2k)}
\).
Pairing \(k\) with \(q-k\), and noting that \(q-2k\equiv q\ (\mathrm{mod}\ 2)\) so \((-1)^{n(q-2k)}=(-1)^{nq}\), the sum of each pair is
\[
(-1)^{kn}\Bigl(\varphi^{n(q-2k)}+(-1)^{n(q-2k)}\varphi^{-n(q-2k)}\Bigr)=(-1)^{kn}L_{n(q-2k)}.
\]
If \(q\) is even, one must also add the contribution of the midpoint mode \(k=q/2\), namely \((-1)^{(q/2)n}\).
The Lucas expression for the primitive orbit count follows directly by substituting into the M\"obius inversion definition. \qed
\end{proof}

\paragraph{Discrete convexity of integer-order pressure and the protocol lower bound on second-order collision}
\begin{corollary}[Discrete convexity and a strict protocol lower bound on \texorpdfstring{$\rho(A_2)$}{rho(A2)}]\label{cor:pom-integer-pressure-discrete-convexity-strict-bound}
Write
\[
\tau_q:=\lim_{m\to\infty}\frac{1}{m}\log S_q(m)
\qquad(q\in\NN),
\]
where \(\tau_0=\log\varphi\) (since \(|X_m|=S_0(m)\asymp \varphi^m\)) and \(\tau_1=\log 2\) (since \(S_1(m)=\sum_x d_m(x)=|\Omega_m|=2^m\)).
Then the integer sequence \((\tau_q)_{q\ge 0}\) satisfies discrete convexity:
\[
\boxed{\;
\tau_{q+1}-\tau_q\ \ge\ \tau_q-\tau_{q-1}\qquad(q\ge 1)\; }.
\]
If moreover the strict convexity regime of Theorem \ref{thm:pom-real-q-pressure-analytic-interpolation} holds (for example, when the fiber multiplicity potential is not cohomologous to a constant), then the above inequality is strict for all \(q\ge 1\).
In particular, in the non-trivial folding (strict convexity) case,
\[
\boxed{\;
\tau_2\ >\ 2\tau_1-\tau_0\ =\ 2\log 2-\log\varphi,
\qquad\text{equivalently}\qquad
\rho\!\left(A_2^{\mathrm{ten}}\right)\ >\ \frac{4}{\varphi}\; }.
\]
\end{corollary}
\begin{proof}
Discrete convexity is the integer-point form of Corollary \ref{cor:pom-crossq-pressure-convexity} (alternatively, one can apply H\"older's inequality directly to each fixed \(m\) to obtain \(S_q(m)^2\le S_{q-1}(m)S_{q+1}(m)\), then take logarithms, divide by \(m\), and pass to the limit).
Strictness follows from the strict convexity condition of Theorem \ref{thm:pom-real-q-pressure-analytic-interpolation}.
Substituting \(q=1\) yields the stated second-order lower bound and spectral radius inequality. \qed
\end{proof}

\paragraph{Removable factor decomposition of the second-order collision spectrum: spectral inversion of the synchronization tail exponent}
\begin{proposition}[Synchronization absorption block and ambiguity tail kernel: removable factors of the collision \texorpdfstring{$\zeta$}{zeta}]\label{prop:pom-collision-zeta-removable-sync-factor}
In the \(q=2\) case of Proposition \ref{prop:pom-momq-tensor-kernel-construction}, decomposing the inverse-code ambiguity state space \(\mathcal S\) into ``synchronized states/unsynchronized states'' (candidate set is a singleton vs.\ non-singleton),
\(A_2^{\mathrm{ten}}\) can be rearranged into a block upper-triangular matrix by permuting the basis:
\[
A_2^{\mathrm{ten}}\ \sim\
\begin{pmatrix}
D & *\\
0 & B
\end{pmatrix},
\]
where \(D\) is the counting kernel of the synchronization absorption block, and \(B\) is the tail spectral kernel of the unsynchronized ambiguity subsystem.
Hence the determinant factorization holds:
\[
\boxed{\;
\det(I-zA_2^{\mathrm{ten}})
=
\det(I-zD)\,\det(I-zB)\; }.
\]
Therefore, once the synchronization absorption factor \(\det(I-zD)\) can be computed explicitly, the ``synchronization tail'' dominant pole radius \(z_{\mathrm{sync}}:=\rho(B)^{-1}\) can be individually inverted from the residual pole set of the second-order collision spectrum
\(\det(I-zA_2^{\mathrm{ten}})\)
after removing \(\det(I-zD)\).
\end{proposition}
\begin{proof}
Since ``synchronized states'' are preserved under the deterministic inverse-code update (once the candidate set becomes a singleton, it never splits again), the subspace spanned by synchronized states is an invariant subspace, yielding the block upper-triangular decomposition.
The determinant of a block triangular matrix equals the product of the diagonal block determinants, hence the factorization formula follows; the pole inversion conclusion is immediate. \qed
\end{proof}

\paragraph{Quadratic characteristic equation induced by the golden two-state skeleton: fixed-point characterization of \texorpdfstring{$\rho(A_q)$}{rho(Aq)}}
\begin{proposition}[Fixed-point equation and monotone root-finding for the Perron root]\label{prop:pom-golden-grammar-fixed-point-rho-aq}
Under the setting of Proposition \ref{prop:pom-momq-tensor-kernel-construction}, \(A_q^{\mathrm{ten}}\) can be written as a \(2\times 2\) block matrix:
\[
A_q^{\mathrm{ten}}
\;=\;
\begin{pmatrix}
T_0^{\otimes q} & T_1^{\otimes q}\\[2pt]
T_0^{\otimes q} & 0
\end{pmatrix}.
\]
Let \(\lambda_q:=\rho(A_q^{\mathrm{ten}})>0\).
Then \(\lambda_q\) is equivalently the unique positive real number satisfying
\[
\boxed{\;
\rho\!\left(\lambda_q^{-1}T_0^{\otimes q}+\lambda_q^{-2}T_1^{\otimes q}T_0^{\otimes q}\right)=1\; }
\]
Moreover, the function
\[
F_q(\lambda):=\rho\!\left(\lambda^{-1}T_0^{\otimes q}+\lambda^{-2}T_1^{\otimes q}T_0^{\otimes q}\right)
\]
is strictly monotonically decreasing in \(\lambda>0\), so \(\lambda_q\) can be found via bisection or monotone iteration with one-dimensional convergence.
\end{proposition}
\begin{proof}
Writing out the Perron eigenpair \((\lambda_q,(u_0,u_1))\) (with \(u_0,u_1\ge 0\) nonzero) from the eigenvalue equation:
\[
\lambda_q u_0=T_0^{\otimes q}u_0+T_1^{\otimes q}u_1,
\qquad
\lambda_q u_1=T_0^{\otimes q}u_0.
\]
Eliminating \(u_1=\lambda_q^{-1}T_0^{\otimes q}u_0\) yields
\(
u_0=\lambda_q^{-1}T_0^{\otimes q}u_0+\lambda_q^{-2}T_1^{\otimes q}T_0^{\otimes q}u_0
\),
so \(1\) is an eigenvalue of \(\lambda_q^{-1}T_0^{\otimes q}+\lambda_q^{-2}T_1^{\otimes q}T_0^{\otimes q}\); by non\-negativity and irreducibility, its spectral radius must equal \(1\).
The converse direction is analogous: if there exists \(\lambda>0\) for which the spectral radius is \(1\), taking the Perron vector \(u_0\) and setting \(u_1=\lambda^{-1}T_0^{\otimes q}u_0\) constructs a Perron eigenvector of \(A_q^{\mathrm{ten}}\), hence \(\lambda=\rho(A_q^{\mathrm{ten}})\).
\par
Monotonicity follows from the monotonicity of the spectral radius with respect to entries of a nonneg\-ative matrix: as \(\lambda\) increases, both terms \(\lambda^{-1}\) and \(\lambda^{-2}\) strictly decrease, so \(F_q\) is strictly decreasing. \qed
\end{proof}

\paragraph{Spectral formula for R\'enyi entropy rates and discrete concavity of the collision index}
\begin{corollary}[Concavity of the collision index and the R\'enyi entropy rate spectral formula]\label{cor:pom-renyi-rate-spectral-formula-discrete-concavity}
For integer \(q\ge 2\), define the unseparated tail
\[
\Pi_q(m):=\sum_{x\in X_m}w_m(x)^q=\frac{S_q(m)}{2^{mq}},
\]
and let its exponential decay rate (collision index) be
\[
E_q:=\lim_{m\to\infty}-\frac{1}{m}\log \Pi_q(m).
\]
If \(S_q(m)=\Theta(\rho(A_q^{\mathrm{ten}})^m)\) (for example, when \(A_q^{\mathrm{ten}}\) is primitive and the boundary vectors are non-degenerate), then
\[
\boxed{\;
E_q=q\log 2-\log\rho(A_q^{\mathrm{ten}})\;=\;h_q\;=\;-P(q)\; }
\]
(notation as in Definition \ref{def:pom-moment-pressure}).
Moreover, the integer sequence \((E_q)_{q\ge 2}\) satisfies discrete concavity in \(q\):
\[
\boxed{\;
E_{q+1}-E_q\ \le\ E_q-E_{q-1}\qquad(q\ge 3)\; }.
\]
The corresponding R\'enyi entropy rate (standard normalization, denoted \(h_q^{\mathrm R}\)) satisfies the spectral closed form
\[
\boxed{\;
h_q^{\mathrm R}
:=\lim_{m\to\infty}\frac{1}{m}\cdot\frac{1}{1-q}\log\sum_{x\in X_m}w_m(x)^q
\;=\;\frac{E_q}{q-1}
\;=\;\frac{q\log 2-\log\rho(A_q^{\mathrm{ten}})}{q-1}\; }.
\]
\end{corollary}
\begin{proof}
The spectral closed form follows directly from \(\Pi_q(m)=S_q(m)/2^{mq}\) and \(\tau(q)=\log\rho(A_q^{\mathrm{ten}})\), and is consistent with the definition \(h_q=\log(2^q/r_q)\) (Definition \ref{def:pom-moment-pressure}).
Discrete concavity follows from the discrete convexity of \(\tau_q\) (the non-strict version of Corollary \ref{cor:pom-integer-pressure-discrete-convexity-strict-bound}) via the identity transformation \(E_q=q\log2-\tau_q\). \qed
\end{proof}
